\section{Objetivo}
O trabalho teve como objetivo implementar os modos de consistência eventual
e Read Your Writes (RYW), executar benchmarks com o YCSB e comparar os
resultados com o modo strong. A seguir, são apresentadas as implementações,
a configuração dos testes e as tabelas de comparação.

\section{Implementação}
O sistema tem um coordenador, implementado em \texttt{coordenador.py}, e três
réplicas que executam \texttt{replica.py}. Cada réplica mantém um arquivo JSON
independente. O cliente acessa o coordenador por HTTP: usa
\texttt{POST /write} para gravar e \texttt{GET /read} para ler. O coordenador
decide quais réplicas devem receber cada pedido conforme o modo escolhido.

\subsection{Consistência eventual}
No eventual, o coordenador escolhe uma réplica por rodízio: 1, 2, 3 e depois
volta à primeira. Ele grava o valor nessa cópia e, quando recebe a confirmação,
agenda a atualização das outras duas e responde sucesso ao cliente.
Assim, o cliente não precisa esperar que todas as cópias sejam atualizadas.

As leituras também escolhem uma réplica por rodízio, usando o mesmo contador
das escritas. Como consultam apenas essa cópia, podem devolver um valor
antigo antes de a sincronização terminar. Por exemplo: se uma escrita muda
o valor de 50 para 80 na réplica 1, uma leitura da réplica 2 ainda pode
retornar 50. Se a chave ainda não chegou à cópia consultada, pode retornar 404.

\subsection{Sincronização e tempo de propagação}
A sincronização é feita pelo coordenador, por meio de uma fila de prioridade
em memória, implementada com \texttt{heapq}. Cada tarefa guarda a chave,
o valor, a versão e as réplicas que precisam receber a atualização. Uma
thread de replicação retira as tarefas quando chega seu horário e envia
os pedidos HTTP aos destinos, um de cada vez.

O parâmetro \texttt{--replication-delay} define o atraso inicial. Nos testes,
ele foi de \textbf{1 segundo}, contado a partir do agendamento após a tentativa
de escrita inicial. O horário é calculado com \texttt{time.monotonic()}.
Esse segundo indica quando o envio pode começar; ele não garante que as
cópias ficarão iguais nesse tempo. A tarefa ainda pode esperar a fila,
os bloqueios e as gravações em disco. Tarefas cujo horário já chegou são
atendidas em sequência, sem uma nova pausa obrigatória de um segundo entre elas.

Quando um envio falha, somente os destinos que falharam voltam à fila.
O atraso da nova tentativa é o maior valor entre 0,1 segundo e o atraso
configurado; portanto, foi de 1 segundo nos testes. O timeout HTTP entre
coordenador e réplica é de 2 segundos, sem contar a espera na fila e nos
bloqueios. Se a tentativa inicial falhar, o cliente recebe 503, mas o
coordenador agenda a escrita para todas as cópias, pois o valor pode ter sido
gravado antes do timeout. Um erro, nesse caso, não significa cancelamento.

Cada escrita recebe uma versão crescente, calculada pelo maior valor entre
o horário atual em nanossegundos e a versão anterior mais um. A réplica
salva a versão junto do valor e recusa versões antigas. Receber a mesma
versão com o mesmo valor apenas confirma a operação, sem regravar o arquivo;
a mesma versão com outro valor gera conflito. Quando surge uma escrita mais
nova para a chave, as tarefas antigas deixam de ser propagadas. As tentativas
continuam enquanto a tarefa for atual e o coordenador estiver ativo.

A fila não é salva em disco, então pendências podem ser perdidas se o
coordenador parar. As réplicas não trocam dados diretamente entre si e não
há uma comparação periódica de todos os arquivos para reconstruir a fila.

\subsection{Read Your Writes (RYW)}
O RYW reutiliza a fila, os tempos e o controle de versões do eventual.
A diferença está na escolha da réplica: cada cliente fica associado a uma
cópia fixa. Para isso, envia o cabeçalho obrigatório \texttt{X-Client-ID}.
O coordenador calcula SHA-256 desse identificador, interpreta os primeiros
oito bytes como um número inteiro e usa o resto da divisão pelo número de
réplicas para escolher o destino.

A escrita é confirmada nessa réplica e as leituras seguintes do mesmo cliente
também vão para ela. Dessa forma, ele consegue ler sua própria escrita
confirmada, ou um valor mais novo, sem esperar a atualização das outras cópias.
Por exemplo, o cliente A pode ler 80 na sua réplica enquanto o cliente B,
associado a outra, ainda lê 50 antes da propagação.

O cliente precisa manter o mesmo identificador, e a lista e a ordem das
réplicas devem continuar iguais. Sem identificador, o coordenador retorna
400. Se a réplica fixa não responder, retorna erro, sem redirecionar para
uma cópia possivelmente atrasada. Clientes diferentes podem cair na mesma
réplica; por isso, a distribuição de carga não é necessariamente equilibrada.
No binding do YCSB, cada thread envia um identificador UUID estável.

\section{Configuração dos benchmarks}
Foram usados 1.000 registros, com um campo de 100 caracteres por registro,
1.000 operações por execução e oito threads. As chaves foram escolhidas
com distribuição uniforme. Houve dez rodadas de cada combinação entre os
três modos e as três cargas: somente leitura, somente escrita (UPDATE)
e mista, com 50\% de probabilidade para cada operação. Ao todo, foram
\textbf{90 execuções e 90.000 operações}.

O YCSB acessou a API por um binding que converte READ em leitura e INSERT/UPDATE
em escrita. Cada teste começou com uma cópia da mesma base sincronizada e
novos processos. A carga inicial ficou fora da medição, os casos rodaram
um por vez e a ordem dos modos alternou entre rodadas. Não houve aquecimento
separado. Cliente e servidores executaram no mesmo computador, com Windows 11,
Python 3.13.7 e Java 21.0.7.

\section{Comparação dos resultados}
As tabelas mostram médias das dez rodadas, com o mesmo peso para cada rodada.
DP significa desvio-padrão, que indica a variação entre elas. Throughput é a
quantidade de operações por segundo: quanto maior, melhor. Latência é o tempo
de resposta: quanto menor, melhor.

\begin{table}[!ht]
\centering\small\setlength{\tabcolsep}{4pt}
\caption{Throughput: média $\pm$ DP, em operações por segundo.}
\begin{tabular}{llll}
\toprule
Modo & Leitura & Escrita & Mista\\\midrule
strong & 41,79 $\pm$ 7,20 & 7,87 $\pm$ 1,46 & 12,46 $\pm$ 1,46\\
eventual & 498,35 $\pm$ 103,50 & 21,22 $\pm$ 3,89 & 40,50 $\pm$ 3,97\\
ryw & 383,85 $\pm$ 101,94 & 19,96 $\pm$ 3,62 & 43,58 $\pm$ 8,48\\
\bottomrule\end{tabular}
\end{table}

O eventual teve a maior média em leitura e escrita: cerca de 11,93 e
2,70 vezes as taxas do strong. Na carga mista, o RYW teve a maior média,
cerca de 3,50 vezes a do strong e 7,59\% acima do eventual. Como houve
variação entre as rodadas, essa diferença não significa que o RYW sempre
será melhor. O rodízio do eventual e a réplica fixa do RYW podem distribuir
o trabalho de formas diferentes, mas esse efeito não foi medido separadamente.

\textbf{O strong demorou mais porque espera as três réplicas antes de confirmar
a escrita e também consulta as três na leitura.} Os pedidos às cópias são
feitos um após o outro, e um bloqueio global faz as operações do coordenador
esperarem umas pelas outras. Além disso, cada escrita lê e regrava o JSON
inteiro da réplica, com sincronização dos dados em disco. O strong espera
esse trabalho nas três cópias; eventual e RYW esperam apenas uma confirmação
e deixam as outras atualizações em segundo plano.

\begin{table}[!ht]
\centering\small\setlength{\tabcolsep}{5pt}
\caption{Latências em ms: média e média dos p95 das rodadas.}
\begin{tabular}{llrrrr}
\toprule
Carga & Modo & Leitura & p95 leitura & Escrita & p95 escrita\\\midrule
Leitura & strong & 196,90 & 275,67 & --- & ---\\
Leitura & eventual & 16,31 & 34,39 & --- & ---\\
Leitura & ryw & 18,27 & 34,85 & --- & ---\\
Escrita & strong & --- & --- & 1.052,94 & 1.603,38\\
Escrita & eventual & --- & --- & 385,44 & 639,44\\
Escrita & ryw & --- & --- & 409,94 & 683,21\\
Mista & strong & 623,18 & 1.002,90 & 674,19 & 1.082,62\\
Mista & eventual & 52,17 & 130,96 & 325,39 & 606,39\\
Mista & ryw & 45,09 & 112,70 & 316,35 & 593,69\\
\bottomrule\end{tabular}
\end{table}

O eventual apresentou menores tempos médios nas cargas isoladas de leitura
e escrita, enquanto o RYW teve menores tempos médios na carga mista.
A leitura do strong passou de 196,90 ms na carga só de leitura para
623,18 ms na mista. Isso é compatível com seu bloqueio global, que faz
leituras esperarem escritas. O tempo no bloqueio não foi medido separadamente.
As latências são das operações bem-sucedidas. O p95 indica o tempo abaixo
do qual ficaram aproximadamente 95\% das operações de uma rodada; a tabela
mostra a média desses p95, não um p95 único de todas as rodadas.

\begin{table}[!ht]
\centering\small
\caption{Duração média do teste, espera adicional pela sincronização e erros.}
\begin{tabular}{llrrr}
\toprule
Carga & Modo & Duração (s) & Espera (s) & Erros\\\midrule
Leitura & strong & 24,72 & 0,01 & 10\\
Leitura & eventual & 2,17 & 0,01 & 45\\
Leitura & ryw & 2,79 & 0,01 & 23\\
Escrita & strong & 132,38 & 0,13 & 20\\
Escrita & eventual & 48,34 & 26,72 & 11\\
Escrita & ryw & 51,43 & 23,22 & 12\\
Mista & strong & 81,36 & 0,02 & 0\\
Mista & eventual & 24,90 & 17,71 & 0\\
Mista & ryw & 23,62 & 15,57 & 1\\
\bottomrule\end{tabular}
\end{table}

Os modos assíncronos terminaram a carga mais rápido, mas ainda precisaram
sincronizar as cópias. Na escrita, essa espera foi de 26,72 s no eventual e
23,22 s no RYW, mostrando que o atraso inicial de 1 segundo não é o tempo
total de propagação. O strong teve pouca espera posterior porque já aguarda
as confirmações durante as escritas.

A espera foi medida após o YCSB, comparando valores e versões dos arquivos,
com intervalo de 0,25 s entre verificações quando havia diferenças.
Ela não entra no throughput nem mede o atraso de cada escrita. Todas as
90 execuções terminaram com as três cópias iguais. Os erros da tabela são
somados nas dez rodadas: foram 30 no strong, 56 no eventual e 36 no RYW.
O throughput inclui as operações que falharam. Esses erros não indicam,
sozinhos, quantas leituras devolveram dados antigos.

Os testes foram locais, sem simulação planejada de falhas de rede.
Os números vêm das execuções salvas; os hashes do código atual diferem dos
registrados nelas, então não foi possível confirmar que o código disponível
é exatamente a versão medida.

\section{Conclusão}
Eventual e RYW foram implementados com propagação em segundo plano e controle
de versões. O RYW acrescenta uma réplica fixa por cliente para permitir a
leitura das próprias escritas. Nos benchmarks, os dois tiveram mais operações
por segundo e menores tempos de resposta que o strong, mas deixaram parte
da sincronização para depois da resposta ao cliente.
